%Chargement des paquets
\usepackage{amsmath}
\usepackage{amsfonts}
\usepackage{amssymb}
\usepackage{enumerate}
\usepackage{mathtools}
\usepackage{bbm}
\usepackage{xparse, etoolbox}
\usepackage{enumerate}
\usepackage{enumitem}
\usepackage{mathabx}
\usepackage{babel}[french]

%environment
\usepackage{amsthm}
\usepackage{keytheorems}
\usepackage{hyperref} 

%Interface théorème
\newkeytheorem{lem}[
	name = Lemme
]

\newkeytheorem{prop}[
	name = Proposition
]

\newkeytheorem{thm}[
	name = Théorème
]

\newkeytheorem{coro}[
	name = Corollaire
]

\renewcommand*{\proofname}{Démonstration}

\theoremstyle{definition}
\newtheorem{defn}{Définition}

\theoremstyle{definition}
\newtheorem*{nota}{Notation}

\theoremstyle{definition}
\newtheorem*{limi}{Liminaire}

\theoremstyle{remark}
\newtheorem{rmq}{Remarque}

\theoremstyle{plain}
\newtheorem*{fait}{Fait}


%Ensembles de nombres
\newcommand{\N} {\mathbb{N}}
\newcommand{\Ne}{\N^\ast}
\newcommand{\Z} {\mathbb{Z}}
\newcommand{\Q} {\mathbb{Q}}
\newcommand{\R} {\mathbb{R}}
\newcommand{\Rb}{\overline{\mathbb{R}}}
\newcommand{\Rp}{\R_+}
\newcommand{\Rm}{\R_-}
\newcommand{\K} {\mathbb{K}}
\newcommand{\Ket} {\mathbb{K^{\ast}}}
\newcommand{\Cx}{\mathbb{C}}

%Ensembles, tribus
\newcommand{\A}{\mathbb{A}}
\renewcommand{\P}{\mathcal{P}}
\newcommand{\T}{\mathcal{T}}
\newcommand{\F}{\mathcal{F}}
\newcommand{\C} {\mathcal{C}}
\newcommand{\ouv}{\mathcal{O}}
\newcommand{\B}{\mathcal{B}}
\newcommand{\M}{\mathcal{M}}
\newcommand{\etag}{\mathcal{E}}
\newcommand{\etagOp}{\etag^{+}(\Omega)}
\newcommand{\etagO}{\etag(\Omega)}
%%Lp avant quotientage
\newcommand{\Lic}{\mathcal{L}^1}
\newcommand{\Lpc}{\mathcal{L}^p}
\newcommand{\Linfc}{\mathcal{L}^{\infty}}
\newcommand{\Lifull}{\Lic(\Omega, \B, m)}
\newcommand{\Lpfull}{\Lpc(\Omega, \B, m)}
\newcommand{\Linffull}{\Linfc(\Omega, \B, m)}
%%Lp après quotientage
\newcommand{\Li}{L^1}
\newcommand{\Ld}{L^2}
\newcommand{\Lp}{L^p}
\newcommand{\Linf}{L^{\infty}}
\newcommand{\Listd}{\Li(\Omega, m)} 
\newcommand{\Ldstd}{\Ld(\Omega, m)} 
\newcommand{\Lpstd}{\Lp(\Omega, m)} 

%Quelques opérateurs
\DeclareMathOperator{\codim}{codim}
\DeclareMathOperator{\rg}{rg}
\DeclareMathOperator{\tr}{tr}
\DeclareMathOperator{\vect}{Vect}
\DeclareMathOperator{\ima}{Im}
\DeclareMathOperator{\id}{Id}
\DeclareMathOperator{\sign}{sign}
\DeclareMathOperator{\ext}{ext}
\newcommand{\ind}{\mathbbm{1}}
\newcommand*{\spr}[2]{\left(\,#1\,\middle|\,#2\,\right)}
\newcommand{\interior}[1]{%
  {\kern0pt#1}^{\mathrm{o}}%
}
\DeclareMathOperator{\card}{card}
\DeclareMathOperator{\ppcm}{ppcm}

%Commandes de pure flemme
\newcommand{\veps}{\varepsilon}
\newcommand{\intab}{\int_{a}^{b}}
\newcommand{\intR}{\int_{-\infty}^{\infty}}
\newcommand{\intinf}{\int_{- \infty}^{+ \infty}}
\newcommand{\equi}{\Leftrightarrow}
\newcommand{\Kev}{$\K$-espace vectoriel }
\newcommand{\Kevs}{$\K$-espaces vectoriels }
\newcommand{\Rev}{$\R$-espace vectoriel }
\newcommand{\Revs}{$\R$-espaces vectoriels }
\newcommand{\Cev}{$\Cx$-espace vectoriel }
\newcommand{\Cevs}{$\Cx$-espaces vectoriels }
\newcommand{\evn}{(E, \norm{.})}
\newcommand{\interf}[2]{[#1,#2]}
\newcommand{\limb}{\lim\limits_}
\newcommand{\lebn}{\lambda_n}
\newcommand{\pp}{\text{~presque partout}}
\newcommand{\derivp}[2]{\frac{\partial #1}{\partial #2}}
\newcommand{\famiu}{(u_i)_{i \in I}}
\newcommand{\sev}{\text{~s.e.v.}}
\newcommand{\Mnp}{M_{n,p}(\K)}
\newcommand{\Mn}{M_{n}(\K)}
\newcommand{\tq}{\text{~t.q.~}}
\newcommand{\mq}{~montrer~que~}
\newcommand{\et}{\quad \text{et} \quad}
\newcommand{\ou}{\text{~ou~}}
\newcommand{\Kx}{\K[X]}


%Commandes de gars chelou
\renewcommand{\subset}{\subseteq}

%Commande restriction, ty duckduckgo
\def\restr#1#2{\mathchoice
              {\setbox1\hbox{${\displaystyle #1}_{\scriptstyle #2}$}
              \restraux{#1}{#2}}
              {\setbox1\hbox{${\textstyle #1}_{\scriptstyle #2}$}
              \restraux{#1}{#2}}
              {\setbox1\hbox{${\scriptstyle #1}_{\scriptscriptstyle #2}$}
              \restraux{#1}{#2}}
              {\setbox1\hbox{${\scriptscriptstyle #1}_{\scriptscriptstyle #2}$}
              \restraux{#1}{#2}}}
\def\restraux#1#2{{#1\,\smash{\vrule height .8\ht1 depth .85\dp1}}_{\,#2}} 

%Quelques normes
\DeclarePairedDelimiter\abs{\lvert}{\rvert}
\DeclarePairedDelimiter\labs{\bigl\lvert}{\bigr\rvert}
\DeclarePairedDelimiter\norm{\lVert}{\rVert}
\DeclarePairedDelimiterXPP\normi[1]{}\lVert\rVert{_1}{\ifblank{#1}{\:\cdot\:}{#1}}
\DeclarePairedDelimiterXPP\normd[1]{}\lVert\rVert{_2}{\ifblank{#1}{\:\cdot\:}{#1}}
\DeclarePairedDelimiterXPP\normp[1]{}\lVert\rVert{_p}{\ifblank{#1}{\:\cdot\:}{#1}}
\DeclarePairedDelimiterXPP\normq[1]{}\lVert\rVert{_q}{\ifblank{#1}{\:\cdot\:}{#1}}
\DeclarePairedDelimiterXPP\norminf[1]{}\lVert\rVert{_\infty}{\ifblank{#1}{\:\cdot\:}{#1}}

%Bracket en angle
\DeclarePairedDelimiter\angl{\langle}{\rangle}

%Set, parenthèses
\newcommand{\set}[1]{\{ #1 \}}
\newcommand{\bigset}[1]{\bigl\{ #1 \bigr\}}
\newcommand{\Bigset}[1]{\Bigl\{ #1 \Bigr\}}
\newcommand{\bigpar}[1]{\bigl( #1 \bigr)}
\newcommand{\Bigpar}[1]{\Bigl( #1 \Bigr)}

%Opitmisation des opérateurs de comparaison
\DeclarePairedDelimiter\tio{o (}{)}
\DeclarePairedDelimiter\gro{O (}{)}


\pagestyle{fancy}

\fancyhead[C]{Théorèmes de Riesz-Fisher}
\fancyhead[R]{}

\begin{document}

\underline{Sources :} Rudin - Analyse - page 79 puis 82.

$X$ est un espace mesuré de mesure $m$. Par convention, $0 \times \infty = 0$.

\begin{thm}[inégalités de Hölder et Minkowski]
Soient $p, q \in [1, \infty[$ vérifiant $\dfrac1{p} + \dfrac1{q} = 1$ et soient $f, g$ deux mesurables sur $X$ et à valeurs dans $[0, \infty]$.
On a :
\begin{align}
\int_X fg dm \leq \left( \int_X f^p dm \right )^{1/p} \left( \int_X g^q dm \right )^{1/q}
\end{align}
et
\begin{align}
\left ( \int_X (f+g)^p dm \right )^{1/p} \leq \left( \int_X f^p dm \right )^{1/p} +  \left( \int_X g^p dm \right )^{1/^p}
\end{align}
\end{thm}

\begin{proof}
\begin{enumerate}[label=(\alph*)]

\item
On note $A$ et $B$ les deux facteurs du membre de droite de l'inégalité (1). Si $A=0$ alors $f$ est nulle presque partout, 
et par suite $fg$  l'est aussi, et l'inégalité est vérifiée. Si $A>0$ et $B = \infty$, l'inégalité est aussi trivialement vraie. \\
On suppose donc que $0 < A, B < \infty$ et on pose :
\[ F = \dfrac{f}{A} \et G = \dfrac{g}{B} , \]
autrement dit (et sans rigueur, car la notion de norme $p$ ou $q$ n'est pas encore démontrée), 
on \og normalise \fg{} les fonctions $f$ et $g$.
Ceci donne :
\[ \int_X F^p dm = \int_X G^p dm = 1 . \]
Pour $x \in X$ vérifiant $0 < F(x), G(x) < \infty$, il existe deux nombres réels $s, t$ tels que :
\[ F(x) = e^{s/p} \et G(x) = e^{t/q} . \]
La fonction exponentielle étant convexe sur $\Rp$, et comme $\dfrac1{p} + \dfrac1{q} = 1$, on a :
\[ e^{\dfrac{s}{p} + \dfrac{t}{q}} \leq \dfrac1{p} e^s + \dfrac1{q} e^t . \]
Autrement dit, pour tout $x \in X$ :
\[ F(x)G(x) \leq \dfrac1{p} F(x)^p + \dfrac1{q} G(x)^q , \]
d'où, en intégrant cette dernière inégalité, on obtient :
\[ \int_X F(x)G(x) dm \leq \dfrac1{p} + \dfrac1{q} = 1 . \]
Il ne reste plus qu'à remplacer $F$ et $G$ par leurs expressions et on conclut.

\item
Comme précédemment, l'inégalité est évidente si le membre de gauche est nulle, soit $f+g$ nulle presque partout.
De même, l'inégalité est évidente si le membre de droite est infini. 
On suppose donc qu'on n'est dans aucun de ces cas et on écrit :
\[ (f+g)^p = f(f+g)^{p-1} + g(f+g)^{p-1} , \]
puis, en appliquant deux fois l'inégalité de Hölder :
\begin{align*}
\int_X (f+g)^p dm & \leq \left [ \left( \int_X f^p dm \right )^{1/p} + \left( \int_X q^p dm \right )^{1/p} \right ] 
 \left ( \int_X  (f+g)^{(p-1)q} \right )^{1/q} \\
& \leq \left [ \left( \int_X f^p dm \right )^{1/p} + \left( \int_X g^p dm \right )^{1/p} \right ]  \left ( \int_X  (f+g)^p dm \right )^{1/q}  
\end{align*}
car $pq = p + q$. Ainsi, cette inégalité sera équivalente à (2) si $\int_X (f+g)^p dm \in ]0, \infty[$. 
Or $f+g$ n'est pas nulle presque partout par hypothèse et, par convexité de la fonction $x \mapsto x^p$, on a 
\[ \left(\dfrac{f+g}2 \right)^p \leq \left( \dfrac1{2} f \right )^p + \left( \dfrac1{2} g \right )^p , \]
ce qui assure que l'intégrale est finie par hypothèse, donc :
\[ \left ( \int_X (f+g)^p dm \right )^{1/p} \leq \left( \int_X f^p dm \right )^{1/p} + \left( \int_X g^p dm \right )^{1/p} \]
car $1-1/q = 1/p$.
\end{enumerate}
\end{proof}

\begin{defn}
Pour $0< p < \infty$, et pour toute fonction mesurable à valeurs complexes définie sur $X$, on pose :
\[ \normp{f} = \left ( \int_X \abs{f}^p dm \right )^{1/p} \]
et on appelle $\Lp(m)$ l'ensemble de toutes les fonctions à valeurs complexes pour lesquelles :
\[ \normp{f} < \infty . \]
Pour $p = \infty$, la norme $\norminf{f}$ est définie comme la plus petite constante $C$ telle que $\abs{f} \leq C$ presque partout,
et l'espace $\Linf$ est celui des fonctions essentiellement bornées.
\end{defn}

\begin{lem}
Les espaces $\Lp(m)$ pour $1 \leq p \leq \infty$ sont des \Cev.
\end{lem}

\begin{thm}[Riesz-Fisher]
Pour $1 \leq p \leq \infty$ et pour toute mesure positive $m$, $\Lp(m)$ est un espace métrique complet.
\end{thm}

\begin{proof}
On suppose d'abord $1 \leq p < \infty$ et soit $(f_n)_{\Ne}$ une suite de Cauchy dans $\Lp$.
Il existe une sous-suite $(f_{n_i})_{\Ne}, n_1 < n_2 < \dots$ telle que :
\[ \forall i \geq 1, \quad \normp{f_{n_{i+1}} + f_{n_i}} < 2^{-i} . \]
On pose :
\[ g_k = \sum_{i=1}^k \abs*{f_{n_{i+1}} - f_{n_i}} \et g = \sum_{i=1}^{\infty} \abs{f_{n_{i+1}} - f_{n_i}} . \]
D'après l'inégalité de Minkowski, on a $\normp{g_k} < 1$ pour $k \geq 1$, donc en appliquant le lemme de Fatou à la suite positive
$(g_k^p)_{\Ne}$ on obtient :
\[ \int_X \liminf g_k^p dm = \int_X g^p dm \leq 1 , \]
de sorte que $\normp{g} \leq 1$. En particulier, $g(x) < \infty$, presque partout, de sorte que la série :
\[ f_{n_1}(x) + \sum_{i=1}^{k-1} (f_{n_{i+1}}(x) - f_{n_i}(x)) , \]
est absolument convergente pour presque tout $x \in E$. On appelle $f(x)$ sa limite lorsque $x$ est tel qu'elle converge, et posons
$f(x)=0$ sur l'ensemble restant, qui est de mesure nulle. Puisque :
\[ f_{n_1} + \sum_{i=1}^{k-1} (f_{n_{i+1}} - f_{n_i}) = f_{n_k} , \]
on voit que :
\[ f(x) = \lim_{i \to \infty} f_{n_i}(x) \pp . \]
La fonction $f$ est donc limite simple presque partout de la suite extraite $(f_{n_i})_{\Ne}$, il reste à démontrer que cette fonction est la 
limite de la suite $(f_n)_{\Ne}$ au sens de la norme $\normp{\cdot}$. \\
Comme cette suite est de Cauchy, pour $\veps >0$, il existe un rang $N$ à partir duquel $\normp{f_{n+k}-f_n} < \veps$ 
pour $n \geq N$ et $k \geq 0$. En appliquant le lemme de Fatou, on obtient :
\[ \int_X \abs{f-f_n}^p dm = \int_X \liminf_k \abs{f_{n+k}-f_n}^p dm \leq \liminf_k \int_X \abs{f_{n+k}-f_n}^p dm \leq \veps^p . \]
En particulier, cette inégalité démontre que $f \in \Lp(m)$ (car $f-f_n \in \Lp$ et $f = f-f_n + f_n$) 
puis que $\normp{f-f_n} \to 0$ pour $n \to \infty$, ce qui achève la démonstration pour $1 \leq p < \infty$. \newline

Pour $p=\infty$, on exhibe à nouveau une suite $(f_n)_{\Ne}$ de Cauchy dans $\Linf(m)$. 
On note $A_k$ les ensembles (négligeables) sur lesquels $\abs{f_k(x)} > \norminf{f_k}$ 
et $B_{n,m}$ les ensembles (eux aussi négligables) sur lesquels $\abs{f_n-f_m} > \norminf{f_n-f_m}$. 
En notant $E$ la réunion de tous ces ensembles, on a $m(E) = 0$ et la suite $(f_n)_{\Ne}$ converge uniformément vers une fonction
bornée $f$ sur le complémentaire de $E$ (chaque suite $(f_n(x))_{\Ne}$ est convergente par complétude de $\Cx$).
On pose $f(x)=0$ sur $E$, clairement $f \in \Linf(m)$ et, par construction $\norminf{f_n - f} \to f$ quand $n \to \infty$.
\end{proof}

\end{document}